\documentclass[hidelinks,12pt,a4paper]{article}
\usepackage[italian]{babel}
\usepackage[utf8]{inputenc}
\usepackage{fourier}

% To avoid GitHub Action error
\usepackage{hyperref}

% Stop hyphenation
\usepackage[none]{hyphenat}

% Justifying text
\emergencystretch 3em

% Enlarge section & subsection
\usepackage{titlesec}
\titleformat*{\section}{\LARGE\bfseries}
\titleformat*{\subsection}{\Large\bfseries}

% Remove first empty page
\usepackage{atbegshi}
\AtBeginDocument{\AtBeginShipoutNext{\AtBeginShipoutDiscard}}

% License
\usepackage[
type={CC},
modifier={by-nc-nd},
version={4.0},
]{doclicense}

\begin{document}
	\begin{flushleft}
		
		%Enlarge text
		\LARGE
		
		\title{\textbf{Copione per podcast sui Musei Civici}}
		\author{Francesco Rombaldoni}
		\date{}
		
		\maketitle
		
		% Adjust page counter
		\setcounter{page}{1}
		\newpage
		
		\tableofcontents
		\newpage
		
		\section{Viaggi Nei Musei di Pesaro: I Musei Civici}
		
		\subsection{Episodio 1: Benvenuti a Palazzo Mosca e la Maledizione della Medusa}
		\begin{center}
			\textbf{Omamori Group Presenta: "Viaggi nei Musei di Pesaro"}
		\end{center}
		Ben ritrovati, cari ascoltatori! Inizia oggi un nuovo ciclo di episodi dedicati a un altro gioiello della nostra città: i Musei Civici di Palazzo Mosca.\\
		Non appena varcherete la soglia di questo splendido palazzo, vi ritroverete faccia a faccia con un'opera imponente e, a dire il vero, un po' inquietante: la gigantesca testa di Medusa in ceramica realizzata da Ferruccio Mengaroni.\\
		Dovete sapere che quest'opera nasconde una storia a dir poco tragica. Si narra che sia avvolta da una vera e propria maledizione! Durante la sua creazione, Mengaroni usò uno specchio per farsi un autoritratto (esatto, la Medusa ha il volto dell'artista intento a urlare!), ma lo specchio si ruppe all'improvviso. Un presagio funesto, se pensiamo che Perseo sconfisse la Medusa proprio usando il suo scudo come specchio.\\
		Purtroppo, la maledizione si avverò: nel 1925, durante l'allestimento della Biennale di Monza, il pesante tondo in ceramica si sbilanciò. Mengaroni corse per salvare il suo capolavoro, ma ne rimase schiacciato, perdendo la vita.\\
		Una storia da brividi che rende quest'opera ancora più magnetica.\\
		E voi, credete alle maledizioni legate agli oggetti d'arte o pensate sia stata solo una tragica fatalità? Fatecelo sapere nei commenti!\\
		Vi ricordo di seguire la nostra pagina Github "Pomodoro Musei di Pesaro". Un grazie speciale a Francesco Rombaldoni per le musiche e il montaggio.\\
		Nella prossima puntata vi parlerò di un quadro famosissimo che nasconde... un altro quadro al suo interno. A presto!\\
		
		\subsection{Episodio 2: Un quadro nel quadro: La Pala di Giovanni Bellini}
		\begin{center}
			\textbf{Omamori Group Presenta: "Viaggi nei Musei di Pesaro"}
		\end{center}
		Bentornati al nostro viaggio tra le sale di Palazzo Mosca! Oggi ci fermiamo davanti al capolavoro indiscusso del museo: la monumentale "Incoronazione della Vergine", dipinta intorno al 1475 dal maestro veneziano Giovanni Bellini.\\
		Quest'opera, dipinta a olio (una tecnica innovativa per l'epoca, importata dalle Fiandre), è immensa e ricchissima di dettagli. Ma la sua particolarità più affascinante è che è costruita come le scatole cinesi!\\
		Se la osservate bene, noterete che la cornice in legno intagliato si fonde perfettamente con l'architettura dipinta nel quadro stesso. Praticamente, Bellini ha creato l'illusione ottica di "un'opera nell'opera". È un trionfo di prospettiva!\\
		Al centro c'è Gesù che incorona Maria, mentre attorno a loro ci sono quattro Santi. Tra questi spicca San Terenzio, patrono di Pesaro. Curiosamente, qui non è vestito da vecchio vescovo, ma da giovane e fiero soldato romano, con in mano un modellino della città.\\
		E voi, quando guardate un grande dipinto, vi soffermate sull'immagine intera o andate a caccia dei piccoli dettagli nascosti sullo sfondo? Scrivetecelo nei commenti!\\
		Grazie sempre a Francesco Rombaldoni per il prezioso lavoro di sound design.\\
		Ma la storia di questo quadro non finisce qui... nel prossimo episodio scopriremo che un pezzo di quest'opera è stato clamorosamente rubato! Curiosi? Continuate a seguirci!\\
		
		\subsection{Episodio 3: Il furto di Napoleone e il salvataggio di Canova}
		\begin{center}
			\textbf{Omamori Group Presenta: "Viaggi nei Musei di Pesaro"}
		\end{center}
		Come vi avevo anticipato, la maestosa Pala del Bellini nasconde una cicatrice storica. Se guardate la grande cornice, noterete che la parte superiore, chiamata "Cimasa", oggi è una ricostruzione. Ma dov'è finita l'originale?\\
		La risposta ha un nome e un cognome ben noti: Napoleone Bonaparte!\\
		Durante le requisizioni francesi, Napoleone rimase così colpito dalla bellezza di questo capolavoro che decise di staccarne la parte superiore, che raffigurava l'Imbalsamazione di Cristo, per portarla in Francia come "souvenir" e bottino di guerra.\\
		Fortunatamente, anni dopo, l'Italia inviò a Parigi una sorta di "agente segreto" dell'arte per recuperare le opere rubate: il celebre scultore Antonio Canova. Grazie alle sue doti diplomatiche, Canova riuscì a riportare in Italia la Cimasa del Bellini. Tuttavia, non tornò mai a Pesaro: oggi è infatti conservata ai Musei Vaticani a Roma!\\
		Solo in rarissime occasioni, come per una grande mostra nel 2008 alle Scuderie del Quirinale, i due pezzi sono stati temporaneamente ricongiunti.\\
		E voi, quale opera d'arte famosa rubata nella storia vorreste veder tornare al suo posto originale? Raccontatecelo nei commenti!\\
		Un enorme grazie a Francesco Rombaldoni per il montaggio. Nel prossimo episodio lasceremo i dipinti per tuffarci nei "social media" del Rinascimento. Non mancate!\\
		
		\subsection{Episodio 4: I "Social Media" del Rinascimento: Le Belle Donne}
		\begin{center}
			\textbf{Omamori Group Presenta: "Viaggi nei Musei di Pesaro"}
		\end{center}
		Oggi entriamo nella luminosa sala dedicata alla collezione di ceramiche di Domenico Mazza. Tra i tanti capolavori in maiolica colorata, ce n'è una categoria molto divertente: la serie delle "Belle Donne".\\
		Immaginate di essere un giovane nobile del Cinquecento in cerca di moglie, o che vuole fare una dichiarazione d'amore. Non potendo mandare un messaggio su WhatsApp, cosa facevate? Semplice: commissionavate un piatto in ceramica finemente decorato!\\
		Questi piatti, chiamati "Doni Amatori", recavano al centro il ritratto idealizzato della fanciulla corteggiata, spesso accompagnato da un nastro dipinto con il suo nome seguito dall'aggettivo "Bella", ad esempio "Faustina Bella".\\
		Erano oggetti di gran lusso, realizzati a Urbino, Casteldurante (oggi Urbania) o nella stessa Pesaro. Alcuni di questi piatti brillano letteralmente d'oro o di rosso rubino grazie all'antica e segreta tecnica del "lustro", una magia chimica ottenuta in cottura.\\
		Insomma, dichiarazioni d'amore scintillanti ed eterne, da esporre in casa per mostrare a tutti le virtù della propria amata.\\
		E a voi, qual è il regalo romantico o di corteggiamento più strano o particolare che abbiate mai fatto o ricevuto? Scrivetecelo qui sotto!\\
		Grazie al nostro Francesco Rombaldoni per il tappeto sonoro di questo episodio. Nella prossima puntata entreremo nell'intimità della camera da letto di una Marchesa insonne... a presto!\\
		
		\subsection{Episodio 5: L'Orologio per insonni della Marchesa Mosca}
		\begin{center}
			\textbf{Omamori Group Presenta: "Viaggi nei Musei di Pesaro"}
		\end{center}
		Cari ascoltatori, oggi facciamo la conoscenza della padrona di casa: la Marchesa Vittoria Toschi Mosca. È grazie a lei, amante dell'antiquariato, che oggi possiamo visitare gran parte di queste collezioni.\\
		La Marchesa viveva in questo palazzo circondata dall'arte, tanto che alla fine decise di donare tutto alla città per ispirare i giovani artisti locali.\\
		Ma tra i lussuosi arredi della sua collezione c'è un oggetto particolarissimo che la Marchesa teneva proprio nella sua camera da letto: l'Orologio Notturno.\\
		La nobildonna soffriva terribilmente di insonnia. Per evitare di dover accendere candele o chiamare la servitù nel cuore della notte per sapere che ore fossero, utilizzava questo geniale orologio del Seicento. Al suo interno veniva posta una candela accesa che retroilluminava i numeri incisi o traforati, permettendole di leggere l'ora nel buio totale!\\
		Sull'orologio, sorretta da un angioletto, c'è una scritta latina: "Volat irreparabile Tempus" (Il tempo vola inesorabilmente). Un monito un po' ansiogeno per chi non riesce a dormire, non trovate?\\
		E voi, che metodi usate per far passare il tempo quando proprio non riuscite a chiudere occhio la notte? Ditecelo nei commenti!\\
		Un ringraziamento a Francesco Rombaldoni alla regia. Vi aspetto al prossimo episodio, dove parleremo di una profezia che lega ancora oggi i cittadini di Pesaro!\\
		
		\subsection{Episodio 6: La Profezia della Beata Michelina}
		\begin{center}
			\textbf{Omamori Group Presenta: "Viaggi nei Musei di Pesaro"}
		\end{center}
		Se c'è una figura a cui i pesaresi sono storicamente devoti, quella è la Beata Michelina. Nata da una ricca famiglia nel 1300, rimase vedova giovanissima. Colpita dal lutto per la perdita dell'unico figlio, decise di donare tutti i suoi averi ai poveri, diventando una terziaria francescana e fondando una confraternita per assistere i malati.\\
		Ai Musei Civici è conservato un magnifico polittico in legno, in stile tardo gotico, realizzato dal pittore veneziano Jacobello Del Fiore, che la ritrae al centro, in una forte posa plastica.\\
		Ma perché Michelina è così importante per la città? Tutto ruota attorno a un'antica leggenda. Prima di morire, consumata dai digiuni, la Beata giurò di proteggere Pesaro per sempre, pronunciando una profezia diventata celebre: "Proteggerò la mia città. Essa tremerà ma non cadrà".\\
		Ancora oggi, ogni volta che la terra trema per un terremoto, i cittadini ricordano queste rassicuranti parole.\\
		Ci sono leggende urbane o figure di protettori simili nella città in cui vivete? Raccontateci le vostre storie nei commenti in basso!\\
		Come sempre, grazie a Francesco Rombaldoni per le musiche e il montaggio. Nel prossimo episodio indagheremo su un vero e proprio "giallo" della storia dell'arte. Non mancate!\\
		
		\subsection{Episodio 7: Un "Cold Case" dell'Arte: La Morte di Adone}
		\begin{center}
			\textbf{Omamori Group Presenta: "Viaggi nei Musei di Pesaro"}
		\end{center}
		Bentornati a "Viaggi nei Musei di Pesaro"! Oggi ci trasformiamo in detective dell'arte per esaminare un dipinto molto intenso e drammatico: la "Morte di Adone".\\
		Questo quadro rappresenta un vero e proprio "cold case" per gli storici dell'arte. Immaginate la scena: Venere e un cupido piangono disperati sul corpo senza vita del giovane e bellissimo Adone.\\
		Il mistero riguarda l'autore. Nell'Ottocento, quando l'opera entrò nelle collezioni, era attribuita con sicurezza a Giovanni Giacomo Sementi. Pochi decenni dopo, nell'euforia di dare lustro alla collezione, venne addirittura promossa a capolavoro del ben più celebre Guido Reni!\\
		E oggi? Colpo di scena! Nell'ultimo catalogo ufficiale del 1993, gli esperti hanno ritrattato tutto, declassando il quadro a opera di "Autore Anonimo". Un declassamento sorprendente per un dipinto dalla fattura così notevole.\\
		E voi, cosa ne pensate? L'importanza o la bellezza di un'opera d'arte per voi cambia se a dipingerla è stato un maestro famoso o un artista rimasto sconosciuto? Fatemi sapere la vostra opinione nei commenti!\\
		Il suono di questa indagine è curato dal nostro Francesco Rombaldoni. Nel penultimo episodio ci aspetta una scena epica e caotica: preparatevi alla caduta dei giganti!\\
		
		\subsection{Episodio 8: Il fragore della battaglia: La Caduta dei Giganti}
		\begin{center}
			\textbf{Omamori Group Presenta: "Viaggi nei Musei di Pesaro"}
		\end{center}
		Oggi ci troviamo nella "Sala del Mito e Devozione" per ammirare un dipinto che definire dinamico è poco. Sto parlando de "La Caduta dei Giganti" del grandissimo pittore Guido Reni.\\
		Immaginate un film d'azione epico, ma dipinto a olio. Il quadro immortala il momento esatto in cui il Re degli Dei, Giove (o Zeus per i greci), scatena la sua furia per cacciare i ribelli Giganti dall'Olimpo.\\
		La composizione è un groviglio di corpi muscolosi e titanici che precipitano rovinosamente verso il basso. Guardando le pennellate decise e le espressioni di terrore dei Giganti sconfitti, sembra quasi di poter sentire il fragore dei tuoni e il rumore assordante dei loro corpi che si schiantano al suolo.\\
		È un'opera di un impatto visivo straordinario, che incarna perfettamente la potenza e la teatralità della mitologia classica.\\
		E voi siete appassionati di mitologia? Qual è il vostro mito greco o romano preferito? Scrivetelo qui sotto nei commenti!\\
		Ringrazio di cuore Francesco Rombaldoni per le musiche epiche di questa puntata. Restate con noi per l'ultimo episodio, dove vi porterò in un luogo del museo dove non dovrete usare gli occhi... ma le orecchie!\\
		
		\subsection{Episodio 9: Oltre la vista: L'esperienza della Sonosfera}
		\begin{center}
			\textbf{Omamori Group Presenta: "Viaggi nei Musei di Pesaro"}
		\end{center}
		Siamo giunti all'ultimo episodio del nostro viaggio a Palazzo Mosca! Per concludere in bellezza, non vi parlerò di dipinti o ceramiche, ma di un'esperienza ultra-tecnologica e unica al mondo che si trova proprio all'interno del polo museale: la Sonosfera.\\
		Pensata e progettata dal Professore del Conservatorio Rossini Davide Monacchi, la Sonosfera è un anfiteatro sferico buio dove non si va per guardare, ma per ascoltare.\\
		Grazie a un complesso sistema di altoparlanti, vi troverete immersi in una bolla sonora a 360 gradi. Potrete assistere allo spettacolo "Frammenti di Estinzione nell'Orologio Climatico", ascoltando i puri suoni registrati nelle più remote Foreste Equatoriali del mondo, per capire quanto sia vitale difendere la nostra Natura.\\
		Oppure, potrete letteralmente "ascoltare" l'arte di Raffaello! I bozzetti dei suoi celebri affreschi della Scuola di Atene o del Parnaso vengono proiettati mentre le musiche dell'epoca e suoni ricostruiti matematicamente vi faranno sentire fisicamente *dentro* il quadro.\\
		È un'esperienza così immersiva e potente che vi cambierà il modo di concepire il suono (consigliata dai ragazzi delle scuole medie in su!).\\
		Vi piacerebbe entrare nella Sonosfera e isolarvi completamente nei suoni della natura? Ditecelo nei commenti!\\
		Finisce qui il nostro tour dei Musei Civici. Un ringraziamento gigantesco a Francesco Rombaldoni per aver curato suono e montaggio di tutta la serie.\\
		Continuate a seguire la pagina Github "Pomodoro Musei di Pesaro" per nuovi contenuti. Grazie di essere stati con noi, e al prossimo viaggio!\\
		
		\vspace*{\fill}
		% Print license shield
		\doclicenseThis
	\end{flushleft}
\end{document}
